\subsection{Isometries and Unitary Matrices}

\begin{definition} \label{def:isometry-linear-map} \hlt{Isometry and Isometric Isomorphism}\\
Let $\mathcal{V}$ and $\mathcal{W}$ be inner-product spaces and let $T\in\LL(\mathcal{V},\mathcal{W})$.
\begin{enumerate}[label=(\roman*)]
\setlength{\itemsep}{0pt}
\item $T$ is an \hlt{isometry} if $\norm{T\mathbf{v}}=\norm{\mathbf{v}}$ for all $\mathbf{v}\in\mathcal{V}$.
\item $T$ is an \hlt{isometric isomorphism} if $T$ is bijective and isometric.\\
The spaces $\mathcal{V}$ and $\mathcal{W}$ are \hlt{isometrically isomorphic} if such a $T$ exists.
\end{enumerate}
\end{definition}

\begin{proposition} \label{prop:basic-isometry-properties} \hlt{Basic Properties of Isometries}\\
Let $\mathcal{V}$ be an inner-product space and let $S,T\in\LL(\mathcal{V})$ be isometries.
\begin{enumerate}[label=\roman*.]
\setlength{\itemsep}{0pt}
\item $ST$ is an isometry;
\item $T$ is injective;
\item if $\mathcal{V}$ is finite-dimensional, then $T$ is invertible and $T^{-1}$ is an isometry.
\end{enumerate}
\begin{proof}
\begin{enumerate}[label=(\roman*)]
\setlength{\itemsep}{0pt}
\item For every $\mathbf{v}$, $\norm{ST\mathbf{v}}=\norm{T\mathbf{v}}=\norm{\mathbf{v}}$.
\item If $T\mathbf{v}=\mathbf{0}$, then $\norm{\mathbf{v}}=\norm{T\mathbf{v}}=0$, so $\mathbf{v}=\mathbf{0}$. Thus $\ker(T)=\{\mathbf{0}\}$.
\item In finite dimensions, injectivity implies invertibility by Theorem \ref{thm:finite-dimensional-operator-invertible-equivalences}.\\
If $\mathbf{u}=T\mathbf{v}$, then $\norm{T^{-1}\mathbf{u}}=\norm{\mathbf{v}}=\norm{T\mathbf{v}}=\norm{\mathbf{u}}$.
\end{enumerate}
\hlt{Structural Commentary}: Isometries preserve length, so they cannot collapse nonzero vectors.\\
\hlt{Key Insight}: In finite dimensions, length preservation already forces an invertible geometry-preserving map.
\end{proof}
\end{proposition}

\begin{example} \label{ex:l2-shift-isometry-not-surjective} \hlt{An Isometry Need Not Be Onto}\\
Define $T:\ell^2(\mathbb{F})\to\ell^2(\mathbb{F})$ by $T(x_1,x_2,\ldots)=(0,x_1,x_2,\ldots)$.\\
For $\mathbf{x}=(x_j)_{j\geq1}$, $\norm{T\mathbf{x}}_2^2=\sum\limits_{j\geq1}\abs{x_j}^2=\norm{\mathbf{x}}_2^2$, so $T$ is an isometry.\\
Since every vector in $\operatorname{range}(T)$ has first coordinate $0$, $(1,0,0,\ldots)\notin\operatorname{range}(T)$; hence $T$ is not surjective.\\
\hlt{Structural Commentary}: Infinite-dimensional isometries can preserve all lengths while missing whole directions.\\
\hlt{Key Insight}: Injective does not imply surjective outside finite dimensions.
\end{example}

\begin{theorem} \label{thm:isometry-iff-inner-product-preserving} \hlt{Isometries Preserve Inner Products}\\
Let $\mathcal{V}$ and $\mathcal{W}$ be inner-product spaces and let $T\in\LL(\mathcal{V},\mathcal{W})$.\\
Then $T$ is an isometry if and only if $\innerprod[\mathcal{W}]{T\mathbf{u},T\mathbf{v}}=\innerprod[\mathcal{V}]{\mathbf{u},\mathbf{v}}$ for every $\mathbf{u},\mathbf{v}\in\mathcal{V}$.
\begin{proof}
\begin{enumerate}[label=\roman*.]
\setlength{\itemsep}{0pt}
\item[$(\Rightarrow)$] Suppose $T$ is an isometry. Then every norm appearing in the polarization identities of Theorem \ref{thm:polarization-identities} is unchanged after applying $T$. Therefore $\innerprod[\mathcal{W}]{T\mathbf{u},T\mathbf{v}}=\innerprod[\mathcal{V}]{\mathbf{u},\mathbf{v}}$.
\item[$(\Leftarrow)$] Suppose $T$ preserves inner products. Then $\norm{T\mathbf{u}}^2=\innerprod[\mathcal{W}]{T\mathbf{u},T\mathbf{u}}=\innerprod[\mathcal{V}]{\mathbf{u},\mathbf{u}}=\norm{\mathbf{u}}^2$, so $T$ is an isometry.
\end{enumerate}
\hlt{Structural Commentary}: Polarization converts preservation of norms into preservation of inner products.\\
\hlt{Key Insight}: Length preservation for all vectors also preserves angle information.
\end{proof}
\end{theorem}

\begin{remark} \label{rem:isometry-conventions} \hlt{Isometry Conventions}
\begin{enumerate}[label=(\roman*)]
\setlength{\itemsep}{0pt}
\item For linear maps between inner-product spaces, norm preservation and inner-product preservation are equivalent by Theorem \ref{thm:isometry-iff-inner-product-preserving}; hence either convention defines the same isometries.
\item For general bilinear forms, an isometry is often required to be bijective and form-preserving; if the form is singular, preservation alone need not imply injectivity.
\end{enumerate}
\end{remark}

\begin{theorem} \label{thm:finite-dimensional-isometry-iff-adjoint-product-identity} \hlt{Isometries and the Adjoint Product}\\
Let $\mathcal{V}$ be finite-dimensional and let $T\in\LL(\mathcal{V})$. Then $T$ is an isometry if and only if $T^*T=I_{\mathcal{V}}$.
\begin{proof}
\begin{enumerate}[label=\roman*.]
\setlength{\itemsep}{0pt}
\item[$(\Rightarrow)$] Suppose $T$ is an isometry. By Theorem \ref{thm:isometry-iff-inner-product-preserving}, $\innerprod{T^*T\mathbf{u},\mathbf{v}}=\innerprod{T\mathbf{u},T\mathbf{v}}=\innerprod{\mathbf{u},\mathbf{v}}$ for all $\mathbf{u},\mathbf{v}$.\\
Thus $\innerprod{(T^*T-I_{\mathcal{V}})\mathbf{u},\mathbf{v}}=0$ for every $\mathbf{v}$, so $T^*T\mathbf{u}=\mathbf{u}$ for every $\mathbf{u}$.
\item[$(\Leftarrow)$] Suppose $T^*T=I_{\mathcal{V}}$. Then $\norm{T\mathbf{u}}^2 =\innerprod{T\mathbf{u},T\mathbf{u}} =\innerprod{T^*T\mathbf{u},\mathbf{u}} =\norm{\mathbf{u}}^2$ for every $\mathbf{u}$, so $T$ is an isometry.
\end{enumerate}
\hlt{Structural Commentary}: The operator $T^*T$ records inner products after applying $T$.\\
\hlt{Key Insight}: $T$ is an isometry exactly when this Gram operator is the identity.
\end{proof}
\end{theorem}

\begin{definition} \label{def:unitary-operator-and-matrix} \hlt{Unitary and Orthogonal}\\
An operator $U\in\LL(\mathcal{V})$ on a finite-dimensional inner-product space is \hlt{unitary} if $U^*U=UU^*=I_{\mathcal{V}}$.\\
A matrix $U\in\operatorname{Mat}(n,n,\mathbb{F})$ is \hlt{unitary} if $U^*U=I_n$. If $\mathbb{F}=\mathbb{R}$, a unitary matrix is called \hlt{orthogonal}.
\end{definition}

\begin{figure}[H]
  \centering
  \begin{tikzpicture}[
  >=Stealth,
  every node/.style={font=\small},
  circleline/.style={gray!50, line width=0.65pt},
  veca/.style={->, blue!75!black, line width=0.95pt},
  vecb/.style={->, teal!70!black, line width=0.95pt},
  map/.style={->, gray!70, line width=0.8pt}
]
  \coordinate (O1) at (0,0);
  \coordinate (O2) at (5.0,0);
  \coordinate (Ua) at ($(O2)+(0.72,0.72)$);
  \coordinate (Ub) at ($(O2)+(-0.72,0.72)$);
  \draw[circleline] (O1) circle (1.15);
  \draw[circleline] (O2) circle (1.15);

  \draw[veca] (O1) -- (1.0,0) node[below] {$\mathbf{e}_1$};
  \draw[vecb] (O1) -- (0,1.0) node[left] {$\mathbf{e}_2$};
  \draw[veca] (O2) -- (Ua) node[above right] {$U\mathbf{e}_1$};
  \draw[vecb] (O2) -- (Ub) node[above left] {$U\mathbf{e}_2$};

  \draw[gray!55] (0.82,0) arc[start angle=0,end angle=90,radius=0.82];
  \draw[gray!55] ($(O2)+(0.58,0.58)$) arc[start angle=45,end angle=135,radius=0.82];
  \node at (0.55,0.55) {$\perp$};
  \node at ($(O2)+(0,0.78)$) {$\perp$};

  \draw[map] (1.65,0.15) -- (3.35,0.15) node[midway, above] {$U$};
\end{tikzpicture}

  \caption{Unitary invariance: $\norm{U\mathbf{x}}=\norm{\mathbf{x}}$ and $\innerprod{U\mathbf{x},U\mathbf{y}}=\innerprod{\mathbf{x},\mathbf{y}}$.}
  \label{fig:unitary-isometry-geometry}
\end{figure}

\begin{theorem} \label{thm:unitary-matrix-equivalent-conditions} \hlt{Equivalent Conditions for a Unitary Matrix}\\
Let $U=[\mathbf{u}_1\ \cdots\ \mathbf{u}_n]\in\operatorname{Mat}(n,n,\mathbb{F})$. The following are equivalent:
\begin{enumerate}[label=\roman*.]
\setlength{\itemsep}{0pt}
\item $U$ is unitary, i.e. $U^*U=I_n$;
\item the columns $\mathbf{u}_1,\ldots,\mathbf{u}_n$ are orthonormal in $\mathbb{F}^n$;
\item $UU^*=I_n$;
\item $U^*$ is unitary;
\item the rows of $U$ are orthonormal;
\item $U$ is invertible and $U^{-1}=U^*$;
\item $\norm{U\mathbf{x}}_2=\norm{\mathbf{x}}_2$ for every $\mathbf{x}\in\mathbb{F}^n$;
\item $\innerprod[\mathbb{F}^n]{U\mathbf{x},U\mathbf{y}}=\innerprod[\mathbb{F}^n]{\mathbf{x},\mathbf{y}}$ for every $\mathbf{x},\mathbf{y}\in\mathbb{F}^n$.
\end{enumerate}
\begin{proof}
\begin{enumerate}[label=(\roman*).]
\setlength{\itemsep}{0pt}
\item[$(i)$]$\Rightarrow (ii)$. Since $(U^*U)_{ij}=\mathbf{u}_i^*\mathbf{u}_j=\innerprod[\mathbb{F}^n]{\mathbf{u}_j,\mathbf{u}_i}$, equation $U^*U=I_n$ states that columns of $U$ are orthonormal.
\item[$(ii)$]$\Rightarrow (iii)$. Orthonormal columns give $U^*U=I_n$. A left inverse of a square matrix is also a right inverse by Theorem \ref{thm:square-one-sided-inverse-is-two-sided}, so $UU^*=I_n$.
\item[$(iii)$]$\Rightarrow (iv)$. The equation $UU^*=I_n$ is $(U^*)^*U^*=I_n$, so $U^*$ is unitary.
\item[$(iv)$]$\Rightarrow (v)$. Apply $(i)\Rightarrow(ii)$ to $U^*$; the columns of $U^*$ are the conjugates of the rows of $U$.
\item[$(v)$]$\Rightarrow (vi)$. Orthonormal rows give $UU^*=I_n$. Thus $U^*$ is the inverse of $U$, so $U^{-1}=U^*$.
\item[$(vi)$]$\Rightarrow (vii)$. If $U^{-1}=U^*$, then $U^*U=I_n$. Therefore $\norm{U\mathbf{x}}_2^2 =\innerprod[\mathbb{F}^n]{U\mathbf{x},U\mathbf{x}} =\innerprod[\mathbb{F}^n]{U^*U\mathbf{x},\mathbf{x}} =\norm{\mathbf{x}}_2^2$ for every $\mathbf{x}\in\mathbb{F}^n$. Since both norms are nonnegative, $\norm{U\mathbf{x}}_2=\norm{\mathbf{x}}_2$.
\item[$(vii)$]$\Rightarrow (viii)$. Apply Theorem \ref{thm:isometry-iff-inner-product-preserving} to $\mathbf{x}\mapsto U\mathbf{x}$.
\item[$(viii)$]$\Rightarrow (i)$. For every $\mathbf{x},\mathbf{y}\in\mathbb{F}^n$, $\innerprod[\mathbb{F}^n]{(U^*U-I_n)\mathbf{x},\mathbf{y}} =\innerprod[\mathbb{F}^n]{U\mathbf{x},U\mathbf{y}}-\innerprod[\mathbb{F}^n]{\mathbf{x},\mathbf{y}} =0$. Taking $\mathbf{y}=(U^*U-I_n)\mathbf{x}$ shows $(U^*U-I_n)\mathbf{x}=\mathbf{0}$ for every $\mathbf{x}$, so $U^*U=I_n$.
\end{enumerate}
\hlt{Structural Commentary}: A unitary matrix is an isometry written in orthonormal coordinates.\\
\hlt{Key Insight}: Columns, rows, inverses, norms, and inner products are different faces of the same condition.
\end{proof}
\end{theorem}

\begin{definition} \label{def:unitary-and-permutation-similarity} \hlt{Unitary and Permutation Similarity}\\
Square matrices $A,B\in\operatorname{Mat}(n,n,\mathbb{F})$ are \hlt{unitarily similar} if $A=UBU^*$ for some unitary $U$. They are \hlt{orthogonally similar} if $\mathbb{F}=\mathbb{R}$ and $U$ can be chosen real orthogonal.\\
A \hlt{permutation matrix} is a square matrix with exactly one entry equal to $1$ in each row and each column, and all other entries equal to $0$. Matrices are \hlt{permutation similar} if $A=PBP^*$ for some permutation matrix $P$.
\end{definition}

\begin{theorem} \label{thm:orthogonal-projection-matrix-unitary-coordinate-form} \hlt{Unitary Coordinate Form of an Orthogonal Projection}\\
Let $P\in\operatorname{Mat}(n,n,\mathbb{F})$ be an orthogonal projection matrix and let $r=\operatorname{rank}(P)$. Then $P$ is unitarily similar to $I_r\oplus0_{n-r}$. Equivalently, there is a unitary matrix $U$ such that
\begin{align}
U^*PU=
\begin{bmatrix}
I_r & 0\\
0 & 0_{n-r}
\end{bmatrix}. \nonumber
\end{align}
\begin{proof}
Let $\mathcal{U}=\operatorname{col}(P)$. Identifying $P$ with its induced operator on $\mathbb{F}^n$, Theorem \ref{thm:orthogonal-projection-self-adjoint-idempotent} says $P$ is the orthogonal projection onto $\mathcal{U}$, so $P$ is the identity on $\mathcal{U}$ and zero on $\mathcal{U}^{\perp}$. Choose an orthonormal basis $\mathbf{u}_1,\ldots,\mathbf{u}_r$ of $\mathcal{U}$ and extend it to an orthonormal basis $\mathbf{u}_1,\ldots,\mathbf{u}_n$ of $\mathbb{F}^n$. Let $U=[\mathbf{u}_1\ \cdots\ \mathbf{u}_n]$.\\
Then $U$ is unitary. For $j\leq r$, $P\mathbf{u}_j=\mathbf{u}_j$, and for $j>r$, $\mathbf{u}_j\in\mathcal{U}^{\perp}$ so $P\mathbf{u}_j=\mathbf{0}$. Therefore the $j$th column of $U^*PU$ is the coordinate column of $P\mathbf{u}_j$ in this orthonormal basis, giving the displayed block diagonal form.\\
\hlt{Structural Commentary}: Orthogonal projection splits the space into range and orthogonal kernel.\\
\hlt{Key Insight}: In a basis adapted to $\operatorname{col}(P)\oplus\operatorname{col}(P)^{\perp}$, projection keeps the first coordinates and kills the rest.
\end{proof}
\end{theorem}

\begin{figure}[H]
  \centering
  \begin{tikzpicture}[
  x=0.88cm,
  y=0.88cm,
  >=Stealth,
  every node/.style={font=\small},
  sub/.style={draw=blue!60!black, fill=blue!7, rounded corners=3pt, minimum width=2.07cm, minimum height=0.64cm, align=center},
  mat/.style={draw=gray!60, fill=gray!5, rounded corners=2pt, minimum width=1.90cm, minimum height=1.58cm},
  arrow/.style={->, line width=0.85pt, blue!70!black}
]
  \node[sub] (C) at (0,0.72) {$\operatorname{col}(P)$};
  \node[sub, fill=red!6, draw=red!55!black] (N) at (0,-0.72) {$\operatorname{col}(P)^\perp$};

  \node[mat] (B) at (4.1,0) {};
  \draw[gray!55] ($(B.center)+(-0.86,0)$) -- ($(B.center)+(0.86,0)$);
  \draw[gray!55] ($(B.center)+(0,-0.68)$) -- ($(B.center)+(0,0.68)$);
  \node[blue!70!black] at ($(B.center)+(-0.43,0.43)$) {$I_r$};
  \node[red!60!black] at ($(B.center)+(0.43,-0.43)$) {$0_{n-r}$};
  \node[gray!65] at ($(B.center)+(0.43,0.43)$) {$0$};
  \node[gray!65] at ($(B.center)+(-0.43,-0.43)$) {$0$};

  \draw[arrow] (C.east) -- ($(B.west)+(0,0.43)$);
  \draw[arrow, red!60!black] (N.east) -- ($(B.west)+(0,-0.43)$);
\end{tikzpicture}

  \caption{Adapted basis: $U^*PU=I_r\oplus0_{n-r}$, identity on $\operatorname{col}(P)$ and zero on $\operatorname{col}(P)^\perp$.}
  \label{fig:orthogonal-projection-unitary-block}
\end{figure}

\begin{proposition} \label{prop:algebra-of-unitary-matrices} \hlt{Algebra of Unitary Matrices}\\
Let $U,V\in\operatorname{Mat}(n,n,\mathbb{F})$ and $W\in\operatorname{Mat}(m,m,\mathbb{F})$, where $V\oplus W$ denotes the block diagonal matrix with blocks $V$ and $W$.
\begin{enumerate}[label=\roman*.]
\setlength{\itemsep}{0pt}
\item If $U$ and $V$ are unitary, then $UV$ is unitary.
\item $V\oplus W$ is unitary if and only if $V$ and $W$ are unitary.
\end{enumerate}
\begin{proof}
\begin{enumerate}[label=(\roman*)]
\setlength{\itemsep}{0pt}
\item Compute $(UV)^*(UV)=V^*U^*UV=V^*V=I_n$.
\item Since $(V\oplus W)^*(V\oplus W)=V^*V\oplus W^*W$, this product equals $I_{n+m}$ exactly when $V^*V=I_n$ and $W^*W=I_m$.
\end{enumerate}
\hlt{Structural Commentary}: Unitary structure is stable under composition and blockwise orthogonal sums.\\
\hlt{Key Insight}: The equation $U^*U=I$ makes every verification a short adjoint calculation.
\end{proof}
\end{proposition}

\begin{theorem} \label{thm:extend-isometric-columns-to-unitary-matrix} \hlt{Extending Isometric Columns to a Unitary Matrix}\\
If $X\in\operatorname{Mat}(m,n,\mathbb{F})$ and $X^*X=I_n$, then there exists $X'\in\operatorname{Mat}(m,m-n,\mathbb{F})$ such that $[X\ X']$ is unitary.
\begin{proof}
Write $X=[\mathbf{x}_1\ \cdots\ \mathbf{x}_n]$. Since $(X^*X)_{ij}=\innerprod[\mathbb{F}^m]{\mathbf{x}_j,\mathbf{x}_i}$, the condition $X^*X=I_n$ says $\mathbf{x}_1,\ldots,\mathbf{x}_n$ are orthonormal in $\mathbb{F}^m$. By Corollary \ref{cor:orthonormal-list-extends-to-basis}, extend this list to an orthonormal basis $\mathbf{x}_1,\ldots,\mathbf{x}_m$ of $\mathbb{F}^m$.\\
Let $X'=[\mathbf{x}_{n+1}\ \cdots\ \mathbf{x}_m]$. Then the columns of $[X\ X']$ are orthonormal, so $[X\ X']$ is unitary.\\
\hlt{Structural Commentary}: Isometric columns are a partial orthonormal basis.\\
\hlt{Key Insight}: Complete the columns by extending the orthonormal list.
\end{proof}
\end{theorem}

\begin{proposition} \label{prop:permutation-matrices-are-orthogonal} \hlt{Permutation Matrices Are Orthogonal}\\
If $P=[\mathbf{e}_{\sigma(1)}\ \cdots\ \mathbf{e}_{\sigma(n)}]$ is a permutation matrix, then $P^*P=PP^*=I_n$. Hence $P$ is real orthogonal.
\begin{proof}
The columns of $P$ are a rearrangement of the standard orthonormal basis of $\mathbb{R}^n$, hence they are orthonormal. Theorem \ref{thm:unitary-matrix-equivalent-conditions} gives $P^*P=PP^*=I_n$.\\
\hlt{Structural Commentary}: A permutation matrix only reorders the standard basis.\\
\hlt{Key Insight}: Reordering orthonormal vectors preserves orthonormality.
\end{proof}
\end{proposition}
